\documentclass[10pt]{article}
\usepackage{tikz}
\usetikzlibrary{arrows.meta, positioning}

\usepackage{parskip}
\usepackage[margin=1in]{geometry} 
\usepackage{amsmath,amsthm,amssymb, graphicx, multicol, array}
\usepackage{enumitem}
\usepackage{amssymb}
\usepackage{float}
\usepackage{href-ul}
 
\newcommand{\N}{\mathbb{N}}
\newcommand{\Z}{\mathbb{Z}}
 
\newenvironment{problem}[2][Problem]{\begin{trivlist}
\item[\hskip \labelsep {\bfseries #1}\hskip \labelsep {\bfseries #2.}]}{\end{trivlist}}

\begin{document}
 
\title{Trends and Cycles II}
\author{Jakob Sverre Alexandersen\\
GRA4159 Trends, Cycles \& Signals Extraction\\
Lecture 7}
\maketitle

\tableofcontents
\newpage
\section{Recap}
\textbf{What we now know}
\begin{itemize}
    \item For estimation, we prefer working with stationary data and models 
    \item Most economic data is non-stationary 
    \item We can remove the non-stationary part using different approaches 
    \begin{itemize}
        \item Differencing the data, i.e., looking at the (first) derivatives / growth 
        \begin{itemize}
            \item But then miss the level which might end up important 
        \end{itemize}
        \item Subtract trend component, but we aware of spurious cycles 
        \begin{itemize}
            \item Deterministic trend is easy, but is it realistic?
            \item Stochastic trends 
        \end{itemize}
    \end{itemize}
\end{itemize}
But often in economics, we are interested in the trend/cycle per se. Not because of the statistical issues involved with non-stationary data, but because theory is written in terms of trends or cycles. 
\section{Trends and cycles}
In the following we will assume that (the log of) a time series can be decomposed into two components, a trend $(g_t) $ and a cycle $(c_t) $:
\begin{gather*}
    y_t = g_t + c_t 
\end{gather*}

To obtain estimates of $(g_t) $ and $(c_t) $, we will consider several univariate detrending methods. 
\section{Deterministic trend}
The traditional method used to estimate business cycles was to define a smooth (natural) growth path for the economy, which was only perturbed by transitory cyclical fluctuations
\begin{gather*}
    y_t = g_t + c_t \\
    \hat{g}_t = \hat{\alpha}_0 + \hat{\alpha}_1t + \hat{\alpha}_2t^2 + \dots\\
    \hat{c}_t = y_t - \hat{g}_t
\end{gather*}
The assumption of a linear trend (LT) implies $|\alpha_1| > 0 \text{, and }\alpha_2 = 0 $.

For a quadratic trend (QT), one assumes $|\alpha_1| > 0 \text{, and }\alpha_2 > 0 $.
\newpage 
\section{Linear filters}
A filter is a function that transforms a time series into another time series. If we let $g_t $ represent a moving average of observed $y_t $, we can extract the trend component $(g_t) $ and the cyclical component $(c_t) $. That is, let $g_t = G(L)y_t $, or: 
\begin{gather}
    g_t = \sum_{j = -m}^{n}g_jy_{t - j}
\end{gather}
where $m, n \in \mathbb{Z}^+$ and $g_j $ are the weights in the $G(L) $ polynomial: 
\begin{gather}
    G(L) = \sum_{j = -m}^{n} g_jL^j
\end{gather}
where $L $ is defined so that $L^jy_t = y_{t - j} $ for positive and negative values of $j $. We will mostly focus our attention to symmetric moving averages, i.e., those for which weights are s.t. $g_j = g_{-j} $.

Equations (1) and (2) define the trend component. The cyclical component is simply the difference between the observed value of $y $ and this component, i.e.:
\begin{gather*}
    c_t = [1 - G(L)]y_t\equiv C(L)y_t 
\end{gather*}
We call $C(L) \text{ and } G(L) $ linear filters. Often the weights are chosen to add up to one, i.e., $\sum_{j = -m}^{n}c_j = 1 $, so that the level of the series is maintained. 

The 5-point moving average filter displayed in lec 1 is an example of a filter that was used to smooth the series. The weight was 1/5, which was applied to a moving window of five observations. That way we created a cyclical component by filtering the noise, i.e.,b
\begin{gather*}
    c_t = 1/5 \sum_{j = -2}^{2}\epsilon_{t - j} = 1/5(\epsilon_{t - 2} + \epsilon_{t - 1} + \epsilon_{t} + \epsilon_{t + 1} + \epsilon_{t + 2})
\end{gather*}
\textbf{Note how applying this type of filters to the data changes its properties!}

One implication: two series that are totally unrelated can become highly related after applying a trend/cycle decomposition
\newpage 
\section{The Hodrick-Prescott (HP) filter}
A pupular way to extract business cycles. The filter extracts a stochastic trend, $(g_t) $, which for a given value of $\lambda $, moves smoothly over time and is uncorrelated with the cycle. The `trend' can be obtained as the solution to the problem: 
\begin{gather*}
    \min_{[g_t]}\sum_{t = 1}^{T}\left[(y_t - g_t)^2 + \lambda[(g_{t + 1} - g_t) - (g_t - g_{t - 1})]^2\right]
\end{gather*}
$\lambda $ is the smoothing parameter. It ``penalizes'' the acceleration in the growth component. 
\begin{itemize}
    \item If $\lambda $ approaches infinity, the lowest minimum is achieved when the variability in the trend is zero, and the trend is perfectly log linear. 
    \item A small value of $\lambda $ will allow more variation in the trend. If $\lambda = 0 $, there will be no cycle, as the trend will follow the observed series perfectly. 
\end{itemize}
\textbf{Warning!}

The HP filter is often used in economics. Use it with care: 
\begin{itemize}
    \item HP filter produces series with spurious dynamic relations that have no basis in the underlying data-generating process. 
    \item The end-of-sample problem: HP filter produces cyclical series that are close to the observed data at the beginning and at the end of the estimation period. 
    \item How to set $\lambda $? And what does it basically imply. $\lambda = 1600 $ is often used for quarterly macro data, but when estimating $\lambda $ (which can be done) one often get very different values 
    \item Read Hamilton on why it sucks: 

    ``(a) The Hodrick-Prescott (HP) filter introduces spurious dynamic relations that have no basis in the underlying data-generating
    process. (b) Filtered values at the end of the sample are very different from
    those in the middle and are also characterized by spurious dynamics. (c)
    A statistical formalization of the problem typically produces values for the
    smoothing parameter vastly at odds with common practice. (d) There is a
    better alternative. A regression of the variable at date $t$ on the four most
    recent values as of date $t - h$ achieves all the objectives sought by users of
    the HP filter with none of its drawbacks.''
\end{itemize}
\newpage 
\section{Beveridge-Nelson (BN) decomposition}
Definition: The implied trend for a non-stationary time series can be defined as the minimum MSE forecast of the long-run level of the series (minus any drift)

An equivalent formulation is that BN trend is the present level of the series plus the infinite sum of the minimum MSE $h $-period ahead first difference forecasts. 

Let $y_t $ be integrated of first order, so that its first differences, $\Delta y_t $, are stationary and can be accurately forecasted using a stationary, e.g., AR(1) model: 
\begin{gather*}
    (\Delta y_t - \alpha) = \phi(\Delta y_{t - 1} - \alpha) + \epsilon_t
\end{gather*}
The minimum squared error (MSE) $h $-period ahead forecast of the first difference AR is (tbc):
\begin{gather}
    \mathbb{E}[\Delta y_{t + h} - \alpha] = \phi^h(\Delta y_t - \alpha)
\end{gather}
As by definition, the BN trend is the present level of the series plus the infinite su of the minimum MSE $h $-period first difference forecasts, we have:
\begin{gather}
    g_t = y_t + \lim_{H\to\infty}\sum_{h = 1}^{H}\mathbb{E}[\Delta y_{t + h} - \alpha]
\end{gather}
Subsituting the expression for the $h $-period ahead forecast from equation (3) into the definition in (4), we get: 
\begin{gather*}
    g_t = y_t + \lim_{H\to\infty}\sum_{h = 1}^{H}\phi^h(\Delta y_t - \alpha)\\
    = y_t + \frac{\phi}{1 - \phi}(\Delta y_t - \alpha)
\end{gather*}
The BN cycle can easily be derived from the result above using the definition: $y_t = g_t + c_t $. Thus: 
\begin{gather*}
    c_t = \frac{\phi}{1 - \phi}(\Delta y_t - \alpha)
\end{gather*}
\subsection{Evaluating the BN filter}
\begin{itemize}
    \item The advantage of the BN method is that it is an appropriate method to extract cycles when a series is difference-stationary and hence using a linear model may infer spurious cycles. 
    \item It allows the series to contain a unit root that can be highly volatile 
    \item One disadvantage is that it is rather time-consuming, as one has to choose between different ARMA models that may give quite different results. 
    \item Cochrane (1988) pointed out that although different ARMA models may fit the short-run properties of the first differences of an observed series, the forecast functions from these models may differ substantially. Since the trend-cycle decomposition in the BN method relies on the forecast properties of the ARMA models, these models may give very different trend-cycle decompositions. 
\end{itemize}
\newpage 
\section{Hamilton (2018) filter}
Builds on similar idea as the BN decomposition, but perhaps even simpler. Builds on fixed horizon direct forecasting equation rather than iterative infinity approximation as in BN. 

If $y_t $ is integrated of some order $d $, do the regression 
\begin{gather*}
    y_{t + h} = \alpha + \sum_{j = 1}^{P}\beta_jy_{t - j + 1}+ v_{t + h}
\end{gather*}
with $p > d $. Then, as shown in Hamilton (2018), the residuals $\hat{v}_{t + h} $ will be stationary and approximate will the cyclical component for a broad class of underlying processes

Thus, $c_{t + g}\equiv \hat{v}_{t + h} = y_{t + h} - q_t $, from which we can easily obtain the trend as well. 
\subsection{Intuition}
Assume $y_t $ is $I(d = 1) $, s.t. $y_t = y_{t - 1} + u_t \text{, and } p = 1 \text{ and } h = 2 $ for simplicity, then the regression 
\begin{gather*}
    y_{t + 2} = \beta y_1 + v_{t + 2}
\end{gather*}
implies that $\hat{\beta} = 1 $ and $\hat{v}_{t + 2} = u_{t + 2} + u_{t + 1} $ because the forecast of an RW, process $h $ periods into the future
\begin{itemize}
    \item Is in expectation $\mathbb{E}[y_{t + h}| y_t] = y_t $. This gives $\hat{\beta} = 1 $ above 
    \item But has errors $y_{t + h} - \mathbb{E}[y_{t + h}|y_t] = \sum_{j = 1}^{h}u_{t + j} $. In turn, this implies that $\hat{v}_{t + 2} = u_{t + 1} + u_{t + 2} $ above 
\end{itemize}
Clearly, since $u_t $ is Gaussian white noise and stationary, $v_{t + 2} $ is stationary as well. Put differently, the sum of stationary components is stationary. 

Note, if the RW assumption is correct, simply computing $y_t - y_{t - h} = \Delta^hy_t = \left(1 - L^h\right)y_t $ would give the ``same'' as OLS regression 

Key benefits (relative to e.g., HP filter)
\begin{itemize}
    \item As BN decomposition; very data driven
    \item Less likely that any inferred correlation between two filtered variables are due to the filtering method per se. Rather, a property of the data. 
    \item Relative to BN, potentially more control, because $h $ can be chosen by the researcher
\end{itemize}

Implementation: 
\begin{itemize}
    \item How to choose $h $: up to you. But, for macroeconomic data and business cycles, Hamilton recommends $h = 8 $ for quarterly data, i.e., 2-year horizon forecast
    \item How to choose $p $: needs to be bigger than $d $ (which is unknown). Hamilton recommends $p = 4 $
\end{itemize}
\newpage 
\section{Spurious cycles}
Nelson and Kang (1981): If one detrend data that are generated by an RW using a smooth deterministic trend, one can infer spurious cycles in the data. The deterministic trend would over-estimate the persistence of the cycle, whereas the trend component would be under-estimated.

Yule-Slutsky effect: A cycle could be spurious, simply because it is a function of the filters applied to extract the cycle from economic data. I.e., using a summation procedure to smooth a series in addition to differentiating to remove the long-run trend, would generate a spurious cyclical pattern. 

Should check for robustness of results (alternative filters)
\section{Summary}
\begin{enumerate}
    \item When studying economic data we are often interested in capturing the trend and cyclical properties of these data series 
    \item You should know the intuition for and be able to apply: 
    \begin{itemize}
        \item Deterministic trend extraction methods 
        \item Moving average filters 
        \item HP method 
        \item BN decomposition
        \item Hamilton filter 
    \end{itemize}
    \item And also understand the strengths and weaknesses with each of these methods 
    \item Remember, the relationship between variables can be very much influenced by different filtering methods 
    \item It is also very easy to generate suprious cycles if one makes filtering assumptions that are at odds with the underlying data generating process. 
\end{enumerate}
\end{document}